% GENERATED FILE. Edit canonical lessons, then run npm run generate:books.
% canonical-source-hash: fnv1a64:72a53339a1db914c
% canonical-lessons: LA-C529-noverca, LA-C529-vitricus, LA-C529-privignus, LA-C529-filiola, LA-C529-fraternus

\chapter{Step-families}
\label{ch:la-step-families}
\hlchaptermodality{529}

\begin{chapteropening}
\textbf{By the end of this chapter:} \emph{I can talk about step-parents and step-children.}

Complete the last lesson of chapter 529: I can talk about step-parents and step-children.
\end{chapteropening}

\section[noverca, novercae]{noverca, novercae — a stepmother}

\label{lesson:LA-C529-noverca}



\begin{quote}
Before the new one: say the Latin for a cousin, then the Latin for twin.
\end{quote}

\subsection*{You'll want to know: noverca, novercae}
\textbf{noverca, novercae} — \textquotedblleft{}a stepmother\textquotedblright{}.

\begin{quote}
Noverca fīliōs marītī amat.
\end{quote}

\begin{quote}
\emph{The stepmother loves her husband's sons.}
\end{quote}

Feminine, first declension. In Roman stories the stepmother is often cruel, but not always.

\begin{cousinweb}
\textbf{noverca} is usually traced to \emph{novus} \textquotedblleft{}new\textquotedblright{} (Lewis \& Short; EtymDB): the new wife of the house.
\end{cousinweb}

\begin{grammarlens}[title={What you've built}]
One more word for telling, explaining and describing at B1.
\end{grammarlens}

\subsection*{Your turn}
\begin{itemize}
\raggedright
  \item \emph{Say it:} \emph{noverca, novercae}
  \item \emph{Say it:} \emph{noverca, novercae}, once more
  \item \emph{Say it:} a sentence of your own with \emph{noverca, novercae}
  \item \emph{Recall:} say the Latin for a cousin, then the Latin for twin, then say \emph{noverca, novercae} again
\end{itemize}

\begin{tcolorbox}[breakable,colback=teal!4,colframe=teal!35!black,title={Before you move on}]
What does noverca, novercae mean? (\textquotedblleft{}A stepmother\textquotedblright{}.) Say it once more.
\end{tcolorbox}
\section[vitricus, vitricī]{vitricus, vitricī — a stepfather}

\label{lesson:LA-C529-vitricus}



\begin{quote}
Before the new one: say the Latin for twin, then the Latin for a stepmother.
\end{quote}

\subsection*{You'll want to know: vitricus, vitricī}
\textbf{vitricus, vitricī} — \textquotedblleft{}a stepfather\textquotedblright{}.

\begin{quote}
Vitricus puerum ut fīlium suum amat.
\end{quote}

\begin{quote}
\emph{The stepfather loves the boy like his own son.}
\end{quote}

Masculine, second declension.

\begin{cousinweb}
\textbf{vitricus} is a native Latin noun, and its origin is not settled; neither Lewis \& Short nor EtymDB gives a source.
\end{cousinweb}

\begin{grammarlens}[title={What you've built}]
One more word for telling, explaining and describing at B1.
\end{grammarlens}

\subsection*{Your turn}
\begin{itemize}
\raggedright
  \item \emph{Say it:} \emph{vitricus, vitricī}
  \item \emph{Say it:} \emph{vitricus, vitricī}, once more
  \item \emph{Say it:} a sentence of your own with \emph{vitricus, vitricī}
  \item \emph{Recall:} say the Latin for twin, then the Latin for a stepmother, then say \emph{vitricus, vitricī} again
\end{itemize}

\begin{tcolorbox}[breakable,colback=teal!4,colframe=teal!35!black,title={Before you move on}]
What does vitricus, vitricī mean? (\textquotedblleft{}A stepfather\textquotedblright{}.) Say it once more.
\end{tcolorbox}
\section[prīvignus, prīvignī]{prīvignus, prīvignī — a stepson}

\label{lesson:LA-C529-privignus}



\begin{quote}
Before the new one: say the Latin for a stepmother, then the Latin for a stepfather.
\end{quote}

\subsection*{You'll want to know: prīvignus, prīvignī}
\textbf{prīvignus, prīvignī} — \textquotedblleft{}a stepson\textquotedblright{}.

\begin{quote}
Prīvignus novercae epistulam scrībit.
\end{quote}

\begin{quote}
\emph{The stepson writes a letter to his stepmother.}
\end{quote}

Masculine; \emph{prīvigna} — a stepdaughter.

\begin{cousinweb}
\textbf{prīvignus} is built, Lewis \& Short explains, on \emph{prīvus} \textquotedblleft{}separate\textquotedblright{} and the root of \emph{gignō} \textquotedblleft{}to beget\textquotedblright{}: born of another line.
\end{cousinweb}

\begin{grammarlens}[title={What you've built}]
One more word for telling, explaining and describing at B1.
\end{grammarlens}

\subsection*{Your turn}
\begin{itemize}
\raggedright
  \item \emph{Say it:} \emph{prīvignus, prīvignī}
  \item \emph{Say it:} \emph{prīvignus, prīvignī}, once more
  \item \emph{Say it:} a sentence of your own with \emph{prīvignus, prīvignī}
  \item \emph{Recall:} say the Latin for a stepmother, then the Latin for a stepfather, then say \emph{prīvignus, prīvignī} again
\end{itemize}

\begin{tcolorbox}[breakable,colback=teal!4,colframe=teal!35!black,title={Before you move on}]
What does prīvignus, prīvignī mean? (\textquotedblleft{}A stepson\textquotedblright{}.) Say it once more.
\end{tcolorbox}
\section[fīliola, fīliolae]{fīliola, fīliolae — a little daughter}

\label{lesson:LA-C529-filiola}



\begin{quote}
Before the new one: say the Latin for a stepfather, then the Latin for a stepson.
\end{quote}

\subsection*{You'll want to know: fīliola, fīliolae}
\textbf{fīliola, fīliolae} — \textquotedblleft{}a little daughter\textquotedblright{}.

\begin{quote}
Fīliola mea hodiē quīnque annōs nāta est.
\end{quote}

\begin{quote}
\emph{My little daughter is five years old today.}
\end{quote}

Feminine, a diminutive of \emph{fīlia}: a word of affection.

\begin{cousinweb}
\textbf{fīliola} is a diminutive of \emph{fīlia} \textquotedblleft{}daughter\textquotedblright{} (Lewis \& Short).
\end{cousinweb}

\begin{grammarlens}[title={What you've built}]
One more word for telling, explaining and describing at B1.
\end{grammarlens}

\subsection*{Your turn}
\begin{itemize}
\raggedright
  \item \emph{Say it:} \emph{fīliola, fīliolae}
  \item \emph{Say it:} \emph{fīliola, fīliolae}, once more
  \item \emph{Say it:} a sentence of your own with \emph{fīliola, fīliolae}
  \item \emph{Recall:} say the Latin for a stepfather, then the Latin for a stepson, then say \emph{fīliola, fīliolae} again
\end{itemize}

\begin{tcolorbox}[breakable,colback=teal!4,colframe=teal!35!black,title={Before you move on}]
What does fīliola, fīliolae mean? (\textquotedblleft{}A little daughter\textquotedblright{}.) Say it once more.
\end{tcolorbox}
\section[frāternus, frāterna, frāternum]{frāternus, frāterna, frāternum — brotherly}

\label{lesson:LA-C529-fraternus}



\begin{quote}
Before the new one: say the Latin for a stepson, then the Latin for a little daughter.
\end{quote}

\subsection*{You'll want to know: frāternus, frāterna, frāternum}
\textbf{frāternus, frāterna, frāternum} — \textquotedblleft{}brotherly\textquotedblright{}.

\begin{quote}
Amor frāternus semper manet.
\end{quote}

\begin{quote}
\emph{Brotherly love always endures.}
\end{quote}

From \emph{frāter}. English \emph{fraternal}.

\begin{cousinweb}
\textbf{frāternus} is built on \emph{frāter} \textquotedblleft{}brother\textquotedblright{} (EtymDB; Lewis \& Short), which is inherited from a Proto-Indo-European word for \textquotedblleft{}brother\textquotedblright{}; English \emph{brother} is a cousin. English \emph{fraternal} comes from it.
\end{cousinweb}

\begin{grammarlens}[title={What you've built}]
One more word for telling, explaining and describing at B1.
\end{grammarlens}

\subsection*{Your turn}
\begin{itemize}
\raggedright
  \item \emph{Say it:} \emph{frāternus, frāterna, frāternum}
  \item \emph{Say it:} \emph{frāternus, frāterna, frāternum}, once more
  \item \emph{Say it:} a sentence of your own with \emph{frāternus, frāterna, frāternum}
  \item \emph{Recall:} say the Latin for a stepson, then the Latin for a little daughter, then say \emph{frāternus, frāterna, frāternum} again
\end{itemize}

\begin{tcolorbox}[breakable,colback=teal!4,colframe=teal!35!black,title={Before you move on}]
What does frāternus, frāterna, frāternum mean? (\textquotedblleft{}Brotherly\textquotedblright{}.) Say it once more.
\end{tcolorbox}
